% Ported from an earlier report section (one-time conversion; edit here).
\subsection{Latency and cost}\label{jb:sec:latency}

A judge exists to save time, so its own speed matters as much as its accuracy. This section explains
where the time goes, using a separate latency study (\LatNSeq\ sequential calls per configuration,
warm-ups excluded, and a per-request random string, a \term{nonce}, so that no provider or local prompt
cache could serve a repeated prompt).

\subsubsection{Anatomy of one call}

\begin{figure}[tbp]
\centering
% Where the time goes in one call: stacked medians from the latency study.
\begin{tikzpicture}
\begin{axis}[benchmark,
  xbar stacked, bar width=11pt,
  width=0.9\linewidth, height=5.2cm,
  xmin=0, ymin=-0.6, ymax=3.6,
  ytick=data, yticklabels from table={generated/jb/data/latency-cloud.dat}{label},
  xlabel={Median milliseconds per call (cloud, connection reused)},
  xmajorgrids,
  legend style={at={(0.5,1)}, anchor=south, yshift=5pt, legend columns=3, /tikz/every even column/.append style={column sep=8pt}},
  legend cell align=left,
]
\addplot[fill=okgray!55, draw=black!70] table[x=network, y=idx]{generated/jb/data/latency-cloud.dat};
\addlegendentry{network}
\addplot[fill=okblue, draw=black!70, postaction={pattern=north east lines, pattern color=white}]
  table[x=server, y=idx]{generated/jb/data/latency-cloud.dat};
\addlegendentry{provider server time}
\addplot[fill=okorange!80, draw=black!70, postaction={pattern=dots, pattern color=white}]
  table[x=other, y=idx]{generated/jb/data/latency-cloud.dat};
\addlegendentry{client and rest}
\end{axis}
\end{tikzpicture}

\vspace{4pt}
\begin{tikzpicture}
\begin{axis}[benchmark,
  xbar stacked, bar width=11pt,
  width=0.9\linewidth, height=5.8cm,
  xmin=0, ymin=-0.6, ymax=4.6,
  ytick=data, yticklabels from table={generated/jb/data/latency-local.dat}{label},
  xlabel={Median milliseconds per call (local, Ollama, model already loaded)},
  xmajorgrids,
  legend style={at={(0.5,1)}, anchor=south, yshift=5pt, legend columns=3, /tikz/every even column/.append style={column sep=8pt}},
  legend cell align=left,
]
\addplot[fill=okgray!55, draw=black!70] table[x=load, y=idx]{generated/jb/data/latency-local.dat};
\addlegendentry{model load/check}
\addplot[fill=okgreen, draw=black!70, postaction={pattern=north east lines, pattern color=white}]
  table[x=prefill, y=idx]{generated/jb/data/latency-local.dat};
\addlegendentry{prompt prefill}
\addplot[fill=okorange!80, draw=black!70, postaction={pattern=dots, pattern color=white}]
  table[x=other, y=idx]{generated/jb/data/latency-local.dat};
\addlegendentry{decode and rest}
\end{axis}
\end{tikzpicture}

\caption{A judge call's latency is made of distinct components, and they differ sharply between hosted and local judges. Where the time goes in one call (medians of each component). Top: hosted judges, split into
network, provider server time and the rest. Bottom: local models through Ollama, split into model
load check, prompt prefill and the rest (decoding one token plus overhead).}
\label{jb:fig:anatomy}
\howtoread{Each bar is one median call; its segments add up to about the median wall-clock time.
Segments differ by colour \emph{and} by fill pattern. The two panels have different horizontal scales, so compare
bar lengths only within a panel.}
\end{figure}

\Cref{jb:tab:anatomy} gives the full numbers for this part.

\begin{table}[tbp]
\centering\small
\caption{The numbers behind \cref{jb:fig:anatomy}. Components are separate medians, so they need not add
exactly to the median wall time; the residual is shown as ``rest.''}
\label{jb:tab:anatomy}
\begin{tabular*}{\textwidth}{@{\extracolsep{\fill}}l r r r r r@{}}
\toprule
Cloud call (keep-alive, medians, ms) & wall & server & network & client + rest & $n$ \\
\midrule
Jev 1.13 & 150 & 52 & 98 & 0 & 30 \\
GPT-6 Luna & 1{,}135 & 986 & 108 & 41 & 30 \\
GPT-6 Luna, priority tier & 768 & 614 & 118 & 37 & 20 \\
GPT-6.1 Sol & 2{,}488 & 2{,}354 & 91 & 44 & 30 \\
\midrule
Local call via Ollama (medians, ms) & wall & load & prefill & decode + rest & $n$ \\
\midrule
Qwen3 0.6B & 29 & 7 & 13 & 9 & 30 \\
tev1 0.8B & 47 & 8 & 24 & 15 & 30 \\
Qwen3 8B & 131 & 10 & 108 & 14 & 30 \\
tev1 4B & 165 & 12 & 110 & 43 & 30 \\
nimble 9B & 272 & 10 & 235 & 27 & 30 \\
\bottomrule
\end{tabular*}

\end{table}

\paragraph{Hosted judges: it is the provider's server} A Jev call took \LatJevWall\,ms at the median,
of which \LatJevServer\,ms was server time. A Luna call took \LatLunaWall\,ms. Of the
\LatLunaGap\,ms gap between Luna and Jev, \LatLunaServerShare\,\% is OpenAI server time; for Sol the
share is \LatSolServerShare\,\%. The client adds at most \LatClientMax\,ms, and the TCP round trip is
about \LatRttJev\,ms to Jev's host and \LatRttOpenAI\,ms to OpenAI's. Opening a fresh connection
instead of reusing one adds \LatFreshRange\,ms (connection plus TLS handshake). OpenAI's priority tier cut
Luna's median by \LatPrioCut\,\% in one window of \LatPrioN\ requests, at \PriorityMultiplier$\times$ the
price; that window was not interleaved with the default tier, so treat it as indicative.

\paragraph{Local judges: prefill and residency} A local judge reads a prompt of about
\LatPromptTokens\ tokens and writes \LatDecodeTokens\ token, so reading the prompt (prefill) is
\LatPrefillShareRange\,\% of the call. The expensive event is a \term{cold start}: when a model is not
already in memory, the first call took \LatColdRange\,s. On this host Ollama serves one request at a
time, so local throughput stays flat as more requests are in flight, while Jev's grows
(\cref{jb:tab:concurrency}).

\begin{table}[tbp]
\centering\small
\caption{Latency and throughput as requests in flight rise. $^{*}$ GPT-6 Luna at 8 in flight had
\LatLunaConcErrors\ read timeout; its throughput counts successful requests only.}
\label{jb:tab:concurrency}
\begin{tabular*}{\textwidth}{@{\extracolsep{\fill}}l r r r r r r@{}}
\toprule
 & \multicolumn{3}{c}{p50 latency (ms)} & \multicolumn{3}{c}{throughput (requests/s)} \\
\cmidrule(lr){2-4}\cmidrule(l){5-7}
Requests in flight & 1 & 4 & 8 & 1 & 4 & 8 \\
\midrule
Jev 1.13 & 146 & 160 & 172 & 6.5 & 16.9 & 43.0 \\
GPT-6 Luna & 1{,}120 & 1{,}047 & 910 & 0.9 & 2.5 & 6.8$^{*}$ \\
tev1 0.8B & 51 & 172 & 350 & 19.1 & 22.6 & 21.6 \\
Laya base & 37 & 132 & 294 & 20.8 & 29.1 & 23.2 \\
Qwen3 8B & 124 & 475 & 963 & 7.8 & 7.9 & 8.0 \\
nimble 9B & 315 & 1{,}154 & 2{,}350 & 2.8 & 3.4 & 3.3 \\
\bottomrule
\end{tabular*}

\end{table}

\begin{caution}
Local latency is a property of this host at that time, not of the model. Across the reproduction runs
the same local answers came back with medians that moved by tens of percent, and the host's one-minute
load average ranged over \HoldLoadRange\ during the holdout run. Cloud latency depends on the provider's
load. Ollama's prompt cache also made repeated prompts look several times faster in an early probe,
which is why every latency probe here uses a nonce.
\end{caution}

\subsubsection{Cost at scale}

\Cref{jb:tab:cost} turns the holdout numbers into a monthly picture, assuming an illustrative
\ProjPerDay\ decisions per day for \ProjDays\ days (\ProjPerMonthM\ million decisions).

\begin{table}[tbp]
\centering\small
\caption{Projected monthly API cost and wrong automatic actions at \ProjPerDay\ decisions per day,
using holdout rates \prereg. Local judges show \$0 API cost and exclude hardware, power and operations.}
\label{jb:tab:cost}
\begin{tabular*}{\textwidth}{@{\extracolsep{\fill}}l r r r r@{}}
\toprule
Judge & API cost per month & wrong automatic per month & automatic share & p95 (ms) \\
\midrule
GPT-6.1 Sol & \${}\SolProjUSD & \SolProjWrong & \SolHoldCovPct\,\% & \SolHoldPninetyfive \\
GPT-6 Luna & \${}\LunaProjUSD & \LunaProjWrong & \LunaHoldCovPct\,\% & \LunaHoldPninetyfive \\
Jev 1.13 & \${}\JevProjUSD & \JevProjWrong & \JevHoldCovPct\,\% & \JevHoldPninetyfive \\
Laya base & \${}\LayaProjUSD & \LayaProjWrong & \LayaHoldCovPct\,\% & \LayaHoldPninetyfive \\
nimble 9B & \${}\NimbleProjUSD & \NimbleProjWrong & \NimbleHoldCovPct\,\% & \NimbleHoldPninetyfive \\
Qwen3 0.6B & \${}\QwenPtSixProjUSD & \QwenPtSixProjWrong & \QwenPtSixHoldCovPct\,\% & \QwenPtSixHoldPninetyfive \\
Qwen3 4B & \${}\QwenFourProjUSD & \QwenFourProjWrong & \QwenFourHoldCovPct\,\% & \QwenFourHoldPninetyfive \\
Qwen3 8B & \${}\QwenEightProjUSD & \QwenEightProjWrong & \QwenEightHoldCovPct\,\% & \QwenEightHoldPninetyfive \\
tev1 0.8B & \${}\TevPtEightProjUSD & \TevPtEightProjWrong & \TevPtEightHoldCovPct\,\% & \TevPtEightHoldPninetyfive \\
tev1 4B & \${}\TevFourProjUSD & \TevFourProjWrong & \TevFourHoldCovPct\,\% & \TevFourHoldPninetyfive \\
\bottomrule
\end{tabular*}

\end{table}

\begin{caution}
The wrong-action column assumes that real traffic has the \emph{same mix} of decisions as this
benchmark. It does not: the cases were deliberately built to be hard, and \HoldNSideOrInjection\ of
\HoldN\ holdout cases involve side effects or injected instructions. Real traffic is likely to contain
far fewer traps, so read this column as a comparison between judges on a stress test, not as a forecast.
The cost column is more portable, because it depends only on tokens per decision and list prices.
\end{caution}

At this volume Jev would cost about \$\JevProjUSD\ a month, GPT-6 Luna about \$\LunaProjUSD\ and
GPT-6.1 Sol about \$\SolProjUSD. The cost ratio to Jev (\LunaHoldCostX$\times$ for Luna,
\SolHoldCostX$\times$ for Sol) is the quantity the preregistered rule tests.
